% Irodov VI §6.3, 105--117. Printed pp.245--246; PDF pp.254--255.
% For term labels begin with \({}^{2S+1}L_J\), preserving JS-string safety.
\atomictermsproblem{105}{245}
求多重态\({}^{4}P\)及\({}^{5}D\)的所有可能总角动量模。
\exsolution
\({}^{4}P\)有$S=3/2,L=1$，故$J=1/2,3/2,5/2$；
\({}^{5}D\)有$S=2,L=2$，故$J=0,1,2,3,4$。
由$M_J=\hbar\sqrt{J(J+1)}$，
\[
 \boxed{{}^{4}P:\ \sqrt3\hbar/2,\ \sqrt{15}\hbar/2,\ \sqrt{35}\hbar/2},
\]
\[
 \boxed{{}^{5}D:\ 0,\sqrt2\hbar,\sqrt6\hbar,2\sqrt3\hbar,2\sqrt5\hbar}.
\]
\excheck 总角动量模是$\hbar\sqrt{J(J+1)}$，不是$J\hbar$。

\atomictermsproblem{106}{245}
求最大总角动量及相应谱项：
(a) Na价电子主量子数$n=4$；
(b) $1s^2$闭壳层外有$2p$与$3d$各一个电子。
\exsolution
(a)允许最大轨道$l=3$，单电子自旋$s=1/2$，
$J_{\max}=7/2$，
\[
 \boxed{{}^{2}F_{7/2},\quad M_{\max}=\sqrt{63}\hbar/2}.
\]
(b)两电子最大合轨道$L=1+2=3$，
平行自旋$S=1$，最大总$J=4$，
\[
 \boxed{{}^{3}F_4,\quad M_{\max}=\sqrt{20}\hbar}.
\]
\excheck 两未配对电子自旋平行产生三重态，
不能写成不可能的\({}^{2}F_4\)。

\atomictermsproblem{107}{245}
已知原子在F和D态的总量子数$J$
各有五种允许值，求两态的自旋角动量模。
\exsolution
耦合$J$不同取值数为$2\min(L,S)+1$。
F态$L=3$故$S=2$，给
\[
 \boxed{M_S(F)=\hbar\sqrt6}.
\]
D态$L=2$故仅可知$S\ge2$，给
\[
 \boxed{M_S(D)\ge\hbar\sqrt6}.
\]
\excheck D态并不能由五种$J$值唯一确定$S$。

\atomictermsproblem{108}{245}
某原子项多重度为3，
总角动量模为$\hbar\sqrt{20}$，
求可能的轨道量子数$L$。
\exsolution
多重度$2S+1=3$故$S=1$，
而$J(J+1)=20$故$J=4$。
角动量三角规则给
$|L-1|\le4\le L+1$，
\[
 \boxed{L=3,4,5}.
\]
\excheck 自旋量子数 $S=1$，总角动量 $J=4$；分别将 $L=3,4,5$ 代入三角不等式 $|L-S|\le J\le L+S$，三者均满足。它们依次对应 F、G、H 三个轨道项；其他整数轨道量子数均不满足。

\atomictermsproblem{109}{246}
求光谱项\({}^{x}D_2\)、
\({}^{x}P_{3/2}\)、\({}^{x}F_1\)
可能的多重度$x$。
\exsolution
由$x=2S+1$、$|L-S|\le J\le L+S$。
(a) $L=2,J=2$给$S=0,1,2,3,4$，
故$\boxed{x=1,3,5,7,9}$。
(b) $L=1,J=3/2$给$S=1/2,3/2,5/2$，
故$\boxed{x=2,4,6}$。
(c) $L=3,J=1$给$S=2,3,4$，
故$\boxed{x=5,7,9}$。
\excheck $J$为半整数时$S$也必须为半整数。

\atomictermsproblem{110}{246}
一个原子除闭壳层外有三电子，分别处于s、p、d轨道。
原子具有该构型最大总角动量。
求矢量模型中总自旋和总角动量夹角。
\exsolution
三自旋平行给$S=3/2$，
最大合轨道$L=0+1+2=3$，
最大$J=L+S=9/2$。
由$\boldsymbol J=\boldsymbol L+\boldsymbol S$，
\[
 \cos\theta=\frac{J(J+1)+S(S+1)-L(L+1)}
 {2\sqrt{J(J+1)S(S+1)}}.
\]
代入得到
\[
 \boxed{\theta=\arccos\left[
 \frac{99/4+15/4-12}{2\sqrt{(99/4)(15/4)}}\right]
 \simeq31^\circ}.
\]
\excheck 经典平行耦合图不能忽略量子角动量模的
$\sqrt{j(j+1)}$修正，故夹角并非0。

\atomictermsproblem{111}{246}
原子总角动量为$\hbar\sqrt6$，自旋$S=1$；
自旋与总角动量夹角$\theta=73.2^\circ$。
求光谱项。
\exsolution
由$J(J+1)=6$得$J=2$。
利用$\boldsymbol L=\boldsymbol J-\boldsymbol S$，
\[
 L(L+1)=J(J+1)+S(S+1)
 -2\cos\theta\sqrt{J(J+1)S(S+1)}\simeq6.
\]
故$L=2$为D项，多重度$2S+1=3$，
\[
 \boxed{{}^{3}D_2}.
\]
\excheck 要用三矢量夹角关系，不能直接用角度估计整数$L$。

\atomictermsproblem{112}{246}
由一个p电子和一个d电子组成的双电子系统，
求其所有允许的LS谱项。
\exsolution
两轨道耦合$L=1,2,3$；
自旋$1/2$与$1/2$耦合$S=0,1$。
单重态$J=L$；三重态$J=L-1,L,L+1$，
所以
\[
 \boxed{{}^{1}P_1,\ {}^{1}D_2,\ {}^{1}F_3,\
 {}^{3}P_{0,1,2},\ {}^{3}D_{1,2,3},\ {}^{3}F_{2,3,4}}.
\]
\excheck p和d为不等价电子，无等价壳层的额外泡利删项。

\atomictermsproblem{113}{246}
原子处于\({}^{2}P_{3/2}\)态，
再增加一个d电子，求合成体系可能的光谱项。
\exsolution
原有p态可取$L_1=1,S_1=1/2$，
新增d电子$L_2=2,S_2=1/2$，
重新LS耦合得$L=1,2,3,S=0,1$，
\[
 \boxed{{}^{1}P_1,\ {}^{1}D_2,\ {}^{1}F_3,\
 {}^{3}P_{0,1,2},\ {}^{3}D_{1,2,3},\ {}^{3}F_{2,3,4}}.
\]
\excheck 本题与前题的LS耦合项集合相同，
只是耦合次序的物理出发点不同。

\atomictermsproblem{114}{246}
判定以下电偶极跃迁的禁戒性：
(a) \({}^{2}D_{3/2}\to{}^{2}P_{1/2}\)；
(b) \({}^{3}P_1\to{}^{2}S_{1/2}\)；
(c) \({}^{3}F_3\to{}^{3}P_2\)；
(d) \({}^{4}F_{1/2}\to{}^{4}D_{5/2}\)。
\exsolution
选择规则$\Delta S=0$，$\Delta L=0,\pm1$，
$\Delta J=0,\pm1$（$0\leftrightarrow0$除外），
且宇称须改变。
(a)角动量条件允许。(b)多重度变3到2即$\Delta S\ne0$，
禁戒。(c)由F到P有$\Delta L=-2$，禁戒。
(d)按本版原题扫描，\({}^{4}F_{1/2}\)不满足
$L=3,S=3/2$下的$J\ge3/2$，
而所印$\Delta J=2$也禁戒。
故
\[
 \boxed{\text{(a)角动量允许；(b)(c)禁戒；
 (d)所印项本身不成立。}}
\]
\excheck 书末仅指出(b)(c)禁戒，
与(d)的印刷下标不一致；已保留争议而非静默改题。

\atomictermsproblem{115}{246}
求锂原子$3D$态的总简并度及其物理意义。
\exsolution
“$3D$”在此表示$n=3$、$L=2$，非三重态。
锂有一个价电子，$S=1/2$，
故$J=3/2,5/2$，
两子态简并度分别为4和6。
\[
 \boxed{g=4+6=10}.
\]
\excheck 未考虑精细结构时
也等于$(2L+1)(2S+1)=5\times2$。

\atomictermsproblem{116}{246}
求谱项\({}^{2}P\)、\({}^{3}D\)、\({}^{4}F\)
在其最大可能$J$下的简并度。
\exsolution
最大$J=L+S$，每个$J$下有$2J+1$个磁分量。
\[
 \boxed{g({}^{2}P)=2(1+\tfrac12)+1=4},\quad
 \boxed{g({}^{3}D)=2(2+1)+1=7},\quad
 \boxed{g({}^{4}F)=2(3+\tfrac32)+1=10}.
\]
\excheck 与整组未分裂项总简并度不同，
这里仅统计指定的最大$J$子项。

\atomictermsproblem{117}{246}
某光谱项简并度为7，且$L=3S$。
求该项符号。
\exsolution
$g=2J+1=7$给$J=3$。
由$L=3S$及$|L-S|\le J\le L+S$，
得$2S\le3\le4S$，
即$3/4\le S\le3/2$。
而$L=3S$必须为整数、$S$为整或半整数，
唯一可行$S=1,L=3$。
所以
\[
 \boxed{{}^{3}F_3}.
\]
\excheck 应联立三角耦合条件，不能只凭简并度判定$L$。
